\section{附注}\label{sec:notes-on-differential-forms}

\subsection{第 \ref{sec:tensors} 节}\label{subsec:notes-tensors}
\begin{enumerate}[label=(\arabic*),leftmargin=*]
\item 给定有限维向量空间 $V$, 定义其\textbf{对偶空间}为 $V$ 上的线性泛函所成的空间.

\begin{RiemannProposition}\label{prop:existence-and-uniqueness-of-dual-basis}
若 $\{v_1,\ldots,v_n\}$ 是 $V$ 的一组基, 则 $V^*$ 存在唯一一组与 $\{v_1,\ldots,v_n\}$ 对偶的基 $\{T_1,\ldots,T_n\}$, 即满足 $T_i(v_j)=\delta_{ij}$ 的基.
\end{RiemannProposition}

\begin{RiemannProof}
由线性可知, 等式 $T_i(v_j)=\delta_{ij}$ 唯一确定了一组泛函 $T_i\in V^*$. 事实上, 对任意 $v\in V$, 有 $v=\sum_{j=1}^na_jv_j$, 从而
\[
  T_i(v)=\sum_{j=1}^na_jT_i(v_j)=\sum_{j=1}^na_j\delta_{ij}=a_i.
\]
此外, 这组唯一确定的泛函线性无关. 事实上, 若
\[
  T:=\sum_{i=1}^nb_iT_i=0,
\]
则对每个 $j=1,\ldots,n$, 有
\[
  0=T(v_j)=\sum_{i=1}^nb_iT_i(v_j)=b_j.
\]
为证明 $\{T_1,\ldots,T_n\}$ 张成 $V^*$, 任取 $S\in V^*$, 令 $b_i:=S(v_i)$. 然后定义 $T:=\sum_{i=1}^nb_iT_i$, 则对所有 $j=1,\ldots,n$ 都有 $S(v_j)=T(v_j)$. 由于 $\{v_1,\ldots,v_n\}$ 是 $V$ 的一组基, 得 $S=T$.
\end{RiemannProof}

此外, 若 $\{v_1,\ldots,v_n\}$ 是 $V$ 的一组基, $\{T_1,\ldots,T_n\}$ 是它的对偶基, 则对任意 $v=\sum a_jv_j\in V$ 和 $T=\sum b_iT_i\in V^*$, 有
\[
  T(v)=\sum_{j=i}^nb_iT_i(v)=\sum_{i,j=1}^na_jb_iT_i(v_j)
  =\sum_{i,j=1}^na_jb_i\delta_{ij}=\sum_{i=1}^na_ib_i.
\]
若现在考虑 $V^*$ 上的线性泛函 $F$, 即 $(V^*)^*$ 的一个元素, 则对某个固定向量 $v_0\in V$, 有 $F(T)=T(v_0)$. 事实上, 设 $\{v_1,\ldots,v_n\}$ 是 $V$ 的一组基, $\{T_1,\ldots,T_n\}$ 是其对偶基. 那么, 若 $T=\sum_{i=1}^nb_iT_i$, 则 $F(T)=\sum_{i=1}^nb_iF(T_i)$. 将 $F(T_i)$ 的值记为 $a_i$, 得 $F(T)=\sum_{i=1}^na_ib_i=T(v_0)$, 其中 $v_0=\sum_{i=1}^na_iv_i$. 这建立了 $(V^*)^*$ 和 $V$ 之间的一一对应, 因而可以将 $V$ 看作 $V^*$ 上的线性泛函空间. 对 $v\in V$ 和 $T\in V^*$, 记 $v(T)=T(v)$.

\item 由 $V$ 的一组基 $\{v_1,\ldots,v_n\}$ 换为新基 $\{v_1',\ldots,v_n'\}$, 得到一个\textbf{换基矩阵} $S$, 其第 $j$ 列是新基向量 $v_j'$ 在旧基下的坐标向量. 可以写出符号化的矩阵等式
\[
  (v_1',\ldots,v_n')=(v_1,\ldots,v_n)S.
\]
向量 $v\in V$ (即 $V$ 上的逆变 1-张量) 关于旧基和新基的坐标 (列) 向量 $a$ 和 $b$ 的关系为
\[
  b=\begin{pmatrix}b_1\\\vdots\\b_n\end{pmatrix}
  =S^{-1}\begin{pmatrix}a_1\\\vdots\\a_n\end{pmatrix}=S^{-1}a,
\]
因为必有 $(v_1',\ldots,v_n')b=(v_1,\ldots,v_n)a=(v_1',\ldots,v_n')S^{-1}a$. 另一方面, 若 $\{T_1,\ldots,T_n\}$ 和 $\{T_1',\ldots,T_n'\}$ 是 $\{v_1,\ldots,v_n\}$ 和 $\{v_1',\ldots,v_n'\}$ 的对偶基, 则
\[
  \begin{pmatrix}T_1\\\vdots\\T_n\end{pmatrix}(v_1,\ldots,v_n)
  =\begin{pmatrix}T_1'\\\vdots\\T_n'\end{pmatrix}(v_1',\ldots,v_n')=I
\]
(在上面的符号化矩阵乘法中, 每一项都由将余向量作用在向量上得到). 因此,
\[
  \begin{pmatrix}T_1\\\vdots\\T_n\end{pmatrix}(v_1',\ldots,v_n')S^{-1}=I
  \ \Leftrightarrow\
  S^{-1}\begin{pmatrix}T_1\\\vdots\\T_n\end{pmatrix}(v_1',\ldots,v_n')=I,
\]
这意味着
\[
  \begin{pmatrix}T_1'\\\vdots\\T_n'\end{pmatrix}=S^{-1}
  \begin{pmatrix}T_1\\\vdots\\T_n\end{pmatrix}.
\]
1-张量 $T\in V^*$ (即 $V$ 上的协变 1-张量) 关于旧基 $\{T_1,\ldots,T_n\}$ 和新基 $\{T_1',\ldots,T_n'\}$ 的坐标 (行) 向量 $a=(a_1,\ldots,a_n)$ 和 $b=(b_1,\ldots,b_n)$ 的关系为
\[
  a\begin{pmatrix}T_1\\\vdots\\T_n\end{pmatrix}
  =b\begin{pmatrix}T_1'\\\vdots\\T_n'\end{pmatrix}
  \ \Leftrightarrow\
  aS\begin{pmatrix}T_1'\\\vdots\\T_n'\end{pmatrix}
  =b\begin{pmatrix}T_1'\\\vdots\\T_n'\end{pmatrix}
\]
所以 $b=aS$. 注意, $V$ 上协变 1-张量的坐标向量与 $V$ 的基向量的变换方式相同 (即通过矩阵 $S$ 变换), 而 $V$ 上逆变 1-张量的坐标向量通过这个矩阵的逆变换. 这就是\textquotedblleft 协变\textquotedblright 和\textquotedblleft 逆变\textquotedblright 两个术语的由来.
\end{enumerate}

\subsection{第 \ref{sec:integration-on-manifolds} 节}\label{subsec:notes-integration-on-manifolds}
\begin{enumerate}[label=(\arabic*),leftmargin=*]
\item (\emph{换元定理}) 设 $U,V\subset\mathbb{R}^n$ 是开集, $g:U\to V$ 是微分同胚, $f:V\to\mathbb{R}$ 是可积函数. 则
\[
  \int_V f=\int_U(f\circ g)|\det\mathrm{d}g|.
\]
\item 要在流形上定义光滑对象, 通常可以先在坐标邻域上定义它们, 再用\textbf{单位分解}将各部分粘合起来.

{%
\setcounter{RiemannDefinition}{0}%
\renewcommand{\theHRiemannDefinition}{\theHsubsection.\arabic{RiemannDefinition}}%
\begin{RiemannTheorem}\label{thm:existence-of-partitions-of-unity}
设 $M$ 是光滑流形, $\mathcal{V}$ 是 $M$ 的一个开覆盖. 则 $M$ 上存在一族可微函数 $\{\rho_i\}_{i\in I}$, 使得:
\begin{enumerate}[label=(\roman*),leftmargin=*]
\item 对每个点 $p\in M$, 存在 $p$ 的一个邻域 $U$, 使得除有限个 $\rho_i$ 外, 均有 $U\cap\operatorname{supp}\rho_i=\varnothing$;
\item 对每个点 $p\in M$, 有 $\sum_{i\in I}\rho_i(p)=1$;
\item $0\leq\rho_i\leq1$, 且对某个 $V\in\mathcal{V}$, 有 $\operatorname{supp}\rho_i\subset V$.
\end{enumerate}
\end{RiemannTheorem}

\begin{RiemannRemark}\label{rem:partition-of-unity-subordinate-to-cover}
这族光滑函数 $\rho_i$ 称为\textbf{从属于覆盖 $\mathcal{V}$ 的单位分解}.
\end{RiemannRemark}

\begin{RiemannProof}
先假设 $M$ 紧. 对每个点 $p\in M$, 考虑其坐标邻域 $W_p=\varphi_p(U_p)$, 使它包含在 $\mathcal{V}$ 的一个元素 $V_p$ 内, 且 $\varphi_p(0)=p$ 和 $B_3(0)\subset U_p$ (其中 $B_3(0)$ 表示以 0 为中心, 半径为 3 的球). 然后考虑下列 $C^\infty$ 函数 (参见 Fig.~\ref{fig:bump-function-graphs}):
\RGfigure[.65\textwidth]{bump-function-graphs}{函数 $\lambda$ 和 $h$ 的图像}
\[
  \begin{array}{rcl}
    \lambda:\mathbb{R}&\to&\mathbb{R}\\
    x&\mapsto&
    \begin{cases}
      e^{\frac{1}{(x-1)(x-2)}}&\text{若 }1<x<2\\
      0&\text{否则},
    \end{cases}\\[.7ex]
    h:\mathbb{R}&\to&\mathbb{R}\\
    x&\mapsto&\displaystyle\frac{\int_x^2\lambda(t)\,\mathrm{d}t}
      {\int_1^2\lambda(t)\,\mathrm{d}t},\\[1ex]
    \beta:\mathbb{R}^n&\to&\mathbb{R}\\
    x&\mapsto&h(\lVert x\rVert).
  \end{array}
\]
注意, $h$ 是递减函数, 其值满足 $0\leq h(x)\leq1$, 当 $x\geq2$ 时等于零, 当 $x\leq1$ 时等于 1. 因此可考虑由下式定义的\phantomsection\label{idx:bump-function-definition}\textbf{隆起函数} $\gamma_p:M\to[0,1]$:
\[
  \gamma_p(q)=
  \begin{cases}
    \beta(\varphi_p^{-1}(q))&\text{若 }q\in\varphi_p(U_p)\\
    0&\text{否则}.
  \end{cases}
\]
于是 $\operatorname{supp}\gamma_p=\overline{\varphi_p(B_2(0))}\subset\varphi_p(B_3(0))\subset W_p$, 从而包含在覆盖的某个元素 $V_p$ 内. 此外, $\{\varphi_p(B_1(0))\}_{p\in M}$ 是 $M$ 的开覆盖, 因而可取有限子覆盖 $\{\varphi_{p_i}(B_1(0))\}_{i=1}^k$, 使得 $M=\bigcup_{i=1}^k\varphi_{p_i}(B_1(0))$. 最后取函数
\[
  \rho_i=\frac{\gamma_{p_i}}{\sum_{j=1}^k\gamma_{p_j}}.
\]
注意, $\sum_{j=1}^k\gamma_{p_j}(q)\ne0$, 因为 $q$ 必然在某个 $\varphi_{p_i}(B_1(0))$ 内, 从而 $\gamma_i(q)\ne0$. 此外, $0\leq\rho_i\leq1$, $\sum\rho_i=1$, 且 $\operatorname{supp}\rho_i=\operatorname{supp}\gamma_{p_i}\subset V_{p_i}$.

若 $M$ 不紧, 可使用一个\textbf{紧穷竭}, 即 $M$ 的紧子集构成的序列 $\{K_i\}_{i\in\mathbb{N}}$, 满足 $K_i\subset\operatorname{int}K_{i+1}$ 且 $M=\bigcup_{i=1}^\infty K_i$. 于是按如下方式得到单位分解. 族 $\{\varphi_p(B_1(0))\}_{p\in M}$ 覆盖 $K_1$, 因此可取 $K_1$ 的有限子覆盖
\[
  \{\varphi_{p_1}(B_1(0)),\ldots,\varphi_{p_{k_1}}(B_1(0))\}.
\]
通过归纳, 得到有限个点, 使得
\[
  \{\varphi_{p_1^i}(B_1(0)),\ldots,\varphi_{p_{k_i}^i}(B_1(0))\}
\]
覆盖 $K_i\setminus\operatorname{int}K_{i-1}$ (一个紧集). 然后对每个 $i$, 考虑相应的隆起函数
\[
  \gamma_{p_1^i},\ldots,\gamma_{p_{k_i}^i}:M\to[0,1].
\]
注意, 对每个 $q\in K_i\setminus\operatorname{int}K_{i-1}$, 有 $\gamma_{p_1^i}+\cdots+\gamma_{p_{k_i}^i}>0$ (因为总有这些函数中的一个非零). 如同紧的情形, 可选择这些隆起函数, 使 $\operatorname{supp}\gamma_{p_j^i}$ 包含在 $\mathcal{V}$ 的某个元素内. 还可以使 $\operatorname{supp}\gamma_{p_j^i}\subset\int K_{i+1}\setminus K_{i-2}$ (一个开集). 因此, $\{\gamma_{p_j^i}\}_{i\in\mathbb{N},1\leq j\leq k_i}$ 是局部有限的, 也就是说, 给定一点 $p\in M$, 存在 $p$ 的开邻域 $V$, 使这些函数中只有有限个不在 $V$ 上恒为零. 所以和 $\sum_{i=1}^\infty\sum_{j=1}^{k_i}\gamma_{p_j^i}$ 是 $M$ 上的正的可微函数. 最后令
\[
  \rho_j^i=\frac{\gamma_{p_j^i}}{\sum_{i=1}^\infty\sum_{j=1}^{k_i}\gamma_{p_j^i}},
\]
就得到了所求的单位分解 (从属于 $\mathcal{V}$).
\end{RiemannProof}

\begin{RiemannRemark}\label{rem:compact-exhaustions-of-manifolds}
流形上总存在紧穷竭. 事实上, 若 $U$ 是 $\mathbb{R}^n$ 的有界开集, 则令
\[
  K_i=\left\{x\in U\mid\operatorname{dist}(x,\partial U)\geq\frac{1}{n}\right\}
\]
就容易为 $U$ 构造一个紧穷竭 $\{K_i\}_{i\in\mathbb{N}}$. 若 $M$ 是可微流形, 总可取可数图册 $\mathcal{A}=\{(U_j,\varphi_j)\}_{j\in\mathbb{N}}$, 使每个 $U_j$ 都是有界开集, 从而都有紧穷竭 $\{K_i^j\}_{i\in\mathbb{N}}$. 因此
\[
  \left\{\bigcup_{i+j=l}\varphi_j\left(K_i^j\right)\right\}_{l\in\mathbb{N}}
\]
是 $M$ 的一个紧穷竭.
\end{RiemannRemark}
}
\end{enumerate}

\subsection{第 \ref{sec:stokes-theorem} 节}\label{subsec:notes-stokes-theorem}

(\emph{Fubini 定理}) 设 $A\subset\mathbb{R}^n$ 和 $B\subset\mathbb{R}^m$ 是紧区间, $f:A\times B\to\mathbb{R}$ 为连续函数. 则
\begin{align*}
  \int_{A\times B}f
  &=\int_A\left(\int_B f(x,y)\,\mathrm{d}y^1\cdots\mathrm{d}y^m\right)
    \mathrm{d}x^1\cdots\mathrm{d}x^n\\
  &=\int_B\left(\int_A f(x,y)\,\mathrm{d}x^1\cdots\mathrm{d}x^n\right)
    \mathrm{d}y^1\cdots\mathrm{d}y^m.
\end{align*}

\subsection{文献附注}\label{subsec:differential-forms-literature}

本章的材料在大多数微分几何著作中都可以找到 (例如 [\RGref{bib:ch02-Boo03}, \RGref{bib:ch02-GHL04}]). \RGcite{dC94} 是一部专门论述微分形式及其应用的著作. 对 de Rham 上同调的研究导向了一套优美而有力的理论, 其细节例如可见 \RGcite{BT82}.
